% ==============================================================================
% T0 Theory / FFGFT: Document 283 (Chapter Content, EN)
% Title: FFGFT <-> HLV: the bridge and the fork
% Author: Johann Pascher
% Date: 17 June 2026
% Purpose: Rigorous test of the FFGFT<->HLV connection against the three-part
%          criterion. Finding: a shared Z3 circulant along Kruger's C3 axis
%          (criteria 1+2), a fundamental fork at the symmetry order
%          (5-fold/phi vs 3-fold/sqrt2).
% References: Doc. 206, 230, 271, 276, 281, 282.
% Script: ffgft_hlv_c3_bridge_probe.py
% ==============================================================================

\section*{Abstract}
Marcel Kr\"uger's HLV builds space as a 6D$\to$3D icosahedral cut-and-project
quasicrystal. Whether FFGFT and HLV are connected is tested here strictly against
the three-part criterion: (1)~the same structural object from each model's
\emph{own} geometry, (2)~an explicit map between the derivations, (3)~a divergence
point that rules out mere Stufe-0 scaffolding. Finding: a Z$_3$ circulant -- the
type of FFGFT's mass operator -- demonstrably lives inside HLV's own geometry,
along a C$_3$ axis of the icosahedral group (criteria 1+2 met). But the internal
invariant diverges \emph{fundamentally}: HLV's protected structure rigidly gives
$r=2/\sqrt5$ (5-fold, $\varphi$), while FFGFT requires $r=\sqrt2$ (3-fold, Koide
$Q=2/3$). No cut-and-project parameter reaches $\sqrt2$ ($r(\tau)\le1$ for all
$\tau$); $\sqrt2$ appears in HLV only as an unprotected accidental angle. Result:
a \emph{genuine} connection with a \emph{fundamental} fork at the symmetry order --
no Stufe-0 triviality, no identity. The fork sits in \emph{space}: HLV's phase layer
is its spiral-time operator (not $\theta=0$), and at the time/memory level HLV's
$\chi$ kernel and FFGFT's Doc.~282 kernel are non-Markovian kernels of the same
class -- a bridge candidate (§\,8). All geometric statements are reproducible with
\texttt{ffgft\_hlv\_c3\_bridge\_probe.py}.

% ==============================================================================
\section{Occasion and criterion}
% ==============================================================================

The exchange with Marcel Kr\"uger (HLV/HICE) repeatedly raises the question
whether FFGFT and HLV are connected. An honest answer needs a sharp criterion, or
one ends up with pseudo-connections that merely share scaffolding (both use
Planck units, both have a square root somewhere). The criterion has three parts:
\begin{enumerate}
\item \textbf{Same object:} the same structural object derived from each model's
      \emph{own} geometry -- not from shared scale scaffolding.
\item \textbf{Map:} an explicit map between the two derivations.
\item \textbf{Divergence point:} a point where the models separate -- otherwise the
      ``connection'' is either trivial identity or Stufe-0 scaffolding.
\end{enumerate}
This document applies the criterion concretely. It claims no theoretical
equivalence; it locates precisely \emph{where} FFGFT and HLV meet and \emph{where}
they diverge.

% ==============================================================================
\section{How Kr\"uger builds his space}
% ==============================================================================

The construction is read directly from the HLV homology script (Doc.~271/276): a
\textbf{6D$\to$3D icosahedral cut-and-project quasicrystal}. A $\mathbb{Z}^6$
lattice is split by two projectors into a 3D parallel space (physical) and a 3D
perpendicular space (the ``hidden layer''); an acceptance window in perp space
selects which lattice points project into the physical quasicrystal. Kr\"uger's
exact projector:
\[
P(\tau)=\begin{pmatrix}
1 & -1 & 0 & 0 & \tau & \tau\\
\tau & \tau & 1 & -1 & 0 & 0\\
0 & 0 & \tau & \tau & 1 & -1
\end{pmatrix}\quad(\text{columns normalised}),
\]
with $\tau=\varphi$ (golden) in parallel space and $\tau=\sigma=-1/\varphi$ in perp
space. The six columns are the six icosahedral five-fold axes (one per antipodal
pair). Internal symmetry: icosahedral (5-fold); the point set is aperiodic.
Striking for a bridge: the hidden space is \emph{3D} -- the same dimension count
as FFGFT's generation fibre $\mathbb{C}^3$.

% ==============================================================================
\section{The shared object: a Z$_3$ circulant on the C$_3$ axis}
% ==============================================================================

The icosahedral group contains C$_3$ axes. The cyclic coordinate map
$(x,y,z)\mapsto(z,x,y)$ is an order-3 rotation; it splits Kr\"uger's six columns
into two 3-cycles, $\{c_1,c_3,c_5\}$ and $\{c_2,c_4,c_6\}$. The Gram matrix of such
a C$_3$ orbit is -- forced by the C$_3$ symmetry -- an \emph{exact real Z$_3$
circulant}, precisely the type of FFGFT's mass operator
$S=\mathrm{circ}(t_0,t_1,t_2)$ on $\mathbb{C}^3$:
\[
G_{\text{orbit}}=\mathrm{circ}\!\left(1,\;\tfrac{1}{\sqrt5},\;\tfrac{1}{\sqrt5}\right)
\quad(\tau=\varphi,\ \text{normalised}),
\]
with $r=2g_1/g_0=2/\sqrt5\approx0.894$ and phase $\theta=0$ (real). The perp space
($\tau=\sigma$) gives $r=-2/\sqrt5$. Closed form: $\;|r|=2|\tau|/(1+\tau^2)$, and for
$\tau=\varphi$ the $\varphi$ cancels exactly to $2/\sqrt5$
($1+\varphi^2=\sqrt5\,\varphi$).

\begin{keyresult}[Criteria 1 and 2 met]
A Z$_3$ circulant lives \emph{inside} HLV's own icosahedral geometry -- not imposed,
but as the Gram matrix of a C$_3$ orbit. This is the same object as FFGFT's mass
operator (criterion~1). The map is explicit (criterion~2): FFGFT's Z$_3$ generation
symmetry $\leftrightarrow$ a C$_3$ axis of HLV's icosahedral group, with the orbit
$\{c_1,c_3,c_5\}$ as image.
\end{keyresult}

% ==============================================================================
\section{The fork: 5-fold versus 3-fold}
% ==============================================================================

The divergence point (criterion~3) is sharp and, as three checks show,
\emph{fundamental} -- not closable by any parameter.

\textbf{Only $1/\sqrt5$.} All $6\times6$ inner products of Kr\"uger's columns have
\emph{a single} magnitude: $1/\sqrt5$. The 5-fold geometry knows only the angle
$\arccos(1/\sqrt5)\approx63.4^\circ$; the $45^\circ$ angle ($\cos=1/\sqrt2$) that
FFGFT needs does not occur between the vacuum directions. Every protected 3-orbit
therefore gives $r=2/\sqrt5$.

\textbf{No $\tau$ reaches $\sqrt2$.} The closed form $r(\tau)=2\tau/(1+\tau^2)$ is
capped by $1$ for \emph{every} $\tau$ (maximum at $\tau=1$, cubic). FFGFT requires
$r=\sqrt2\approx1.414>1$ -- which is \emph{outside} the reachable range of
Kr\"uger's construction, for any parameter.

\textbf{$\sqrt2$ needs 3-fold.} In a general 3-fold geometry (cyclic basis
$(a,b,c)$, $r=2(ab+bc+ca)/(a^2+b^2+c^2)$) $\sqrt2$ is perfectly reachable, e.g.
$v=(1,1,\sqrt2-2^{1/4})$. FFGFT's Koide point
($\mathrm{circ}(1,\tfrac1{\sqrt2},\tfrac1{\sqrt2})$, eigenvalues $0.29/0.29/2.41$,
positive semidefinite) is thus geometrically legitimate -- it merely needs a
\emph{3-fold} column, not Kr\"uger's 5-fold one.

\begin{center}
\begin{tabular}{lcccl}
\toprule
 & $r$ & $\theta$ & Koide $Q$ & source \\
\midrule
HLV (protected) & $2/\sqrt5=0.894$ & $0$ & $7/15$ & 5-fold / $\varphi$ \\
FFGFT (mass op.) & $\sqrt2=1.414$ & $2/9$ & $2/3$ & 3-fold / Koide \\
\bottomrule
\end{tabular}
\end{center}

% ==============================================================================
\section{Protected versus generic: no loophole}
% ==============================================================================

One objection remains: the HLV quasicrystal has thousands of projected points --
could a C$_3$ orbit of \emph{some} point accidentally hit $\sqrt2$? For $9822$
projected points in the acceptance window the C$_3$-orbit ratios $g_1/g_0$ do fill
the whole continuum $[-\tfrac12,1]$, and $1/\sqrt2$ is among them $150$ times.

\begin{warningboxenv}[A generic hit proves nothing]
That generic projected points sit at \emph{every} angle -- including $1/\sqrt2$ --
holds for any dense 3-fold point set, even a cubic one. This is the Stufe-0
triviality: one picks an \emph{unprotected} direction with the right angle. The
\emph{symmetry-protected} orbits -- Kr\"uger's six vacuum directions -- by contrast
rigidly give only $\pm1/\sqrt5$. $\sqrt2$ is not present in HLV's protected
structure; it appears only as a meaningless accidental angle.
\end{warningboxenv}

% ==============================================================================
\section{HLV is the decompactified variant: mass and time only through the addition}
% ==============================================================================

One layer of the fork sits deeper than the symmetry order and is easy to miss.
HLV's cut-and-project maps into an infinite, aperiodic $\mathbb{R}^3$ -- it is the
\emph{decompactified} variant. FFGFT's space is by contrast \emph{compact}
($T^4$). This is not cosmetic: mass and time are intrinsic in FFGFT because
compactness quantises the modes -- that is the time-mass duality
$\tilde T\cdot m=1$. Without compactness there is no mode quantisation, hence
\emph{no intrinsic mass scale and no time}. That is exactly why HLV's protected
structure is purely geometric: a real Gram metric with $\theta=0$ and dimensionless
angles. The intrinsic \emph{mass} scale is missing because the decompactified
variant has no mode quantisation; it arrives, as already recorded in Doc.~282, only
\emph{through the addition}. But it matters that this $\theta=0$ is the spatial
metric, \emph{not} HLV's time. HLV is not timeless or phaseless -- its time sits in
the native spiral-time operators (§\,8), which the spatial Gram does not capture at
all. The C$_3$ probe saw only the geometry, not the dynamics.

This sharpens criterion~1. The shared Z$_3$ circulant is the same only
\emph{algebraically} (the same cyclic form). Physically it differs: in HLV
(decompactified) a \emph{spatial metric} without mass and time, in FFGFT (compact)
the \emph{mass operator} with mass and time intrinsic. The bridge carries the form,
not the content.

\begin{keyresult}[The fork has three layers]
\begin{center}
\begin{tabular}{lll}
\toprule
layer & HLV & FFGFT \\
\midrule
compactness & decompactified (metric) & compact ($\tilde T\!\cdot\! m=1$) \\
symmetry order & 5-fold ($\varphi$, $1/\sqrt5$) & 3-fold ($\sqrt2$) \\
phase & spiral-time: skew-adj.\ $R$ (dyn., §\,8) & scalar, $\theta=2/9$ \\
\bottomrule
\end{tabular}
\end{center}
The complex holonomy ($\theta=2/9$, FFGFT) demanded above is the \emph{operational
consequence} of compactness: there mass and phase arise from the quantised modes.
The HLV column of the phase row is subtler than a bare $\theta=0$ -- that was the
\emph{spatial} Gram metric; HLV's own time-phase sits in the skew-adjoint
spiral-time operator $R$ (§\,8). A full recompactification of HLV must
\emph{co-integrate} these operators, not merely periodise the space -- and a compact
HLV is moreover no longer a quasicrystal.
\end{keyresult}

% ==============================================================================
\section{A compactified variant of HLV: the rational approximant}
% ==============================================================================

Can this recompactification be carried out concretely, to make the comparison
fair? Yes -- with an established notion from quasicrystal physics, the
\emph{rational approximant}. One replaces the irrational $\varphi$ in the projector
by rational Fibonacci ratios $\tau_n=F_{n+1}/F_n$
$(1,2,\tfrac32,\tfrac53,\tfrac85,\dots\to\varphi)$. For rational $\tau=p/q$ the
physical projector, multiplied by $q$, becomes \emph{integer} -- the projection of
$\mathbb{Z}^6$ lands on a scaled integer lattice $(1/q)\mathbb{Z}^3$, i.e. on a
\emph{periodic crystal} on $T^3$. For irrational $\varphi$ no such factor exists;
the point set stays dense and aperiodic. That is the precise sense of
``compactified HLV'' -- and, as §\,6 already noted, it is no longer a quasicrystal
but its periodic crystal approximant.

The approximant is periodic and therefore has discrete modes -- a \emph{mass scale
becomes possible}. This closes the deepest layer of the fork (compactness): HLV is
now in the same regime as FFGFT. The C$_3$ ratio converges:

\begin{center}
\begin{tabular}{lcc}
\toprule
$\tau=F_{n+1}/F_n$ & $r$ (C$_3$) & Koide $Q$ \\
\midrule
$2/1$ & $0.800$ & $0.440$ \\
$5/3$ & $0.882$ & $0.463$ \\
$13/8$ & $0.893$ & $0.466$ \\
$\to\varphi$ & $0.894=2/\sqrt5$ & $7/15$ \\
\bottomrule
\end{tabular}
\end{center}

But the approximant inherits HLV's 5-fold symmetry: $r\to2/\sqrt5$, Koide
$\to7/15$; the \emph{spatial} metric stays real ($\theta=0$). So the symmetry order
keeps diverging. The phase layer, however, is not settled by ``$\theta=0$'': the
purely spatial approximant only recompactifies space; HLV's time-phase sits in the
spiral-time operators it still has to integrate (§\,8).

\begin{keyresult}[The fork survives the fair comparison]
With the approximant the comparison is \emph{compact vs compact}, and thus fair. It
confirms the fork: \emph{even compactified}, HLV gives Koide $7/15$, not FFGFT's
$2/3$. The divergence was therefore \emph{not a compact/decompact artefact} -- it
sits genuinely in the 5- vs 3-fold symmetry. One can compactify HLV; the finding
holds.
\end{keyresult}

% ==============================================================================
\section{What compactification must integrate: HLV's spiral-time}
% ==============================================================================

The approximant of §\,7 recompactifies HLV's \emph{space}. But HLV is not purely
spatial: its defining dynamics are the native 6D \emph{spiral-time} operators
\begin{equation}
\Psi(t)=t+i\,\phi(t)+j\,\chi(t),
\end{equation}
with $\phi(t)$ a phase/coherence channel and $\chi(t)$ a memory/hysteresis channel
(Krüger, Zenodo 20668125 and 20512650). A compactified variant that is to
\emph{remain} HLV must integrate these operators -- the purely spatial approximant
does not yet do so. This corrects the phase layer of the fork.

Krüger does not apply $\Psi$ directly (that would give hypercomplex fields) but via
a \emph{conservative translation} into two real operators:
\begin{center}
\begin{tabular}{lll}
\toprule
channel & becomes & requirement \\
\midrule
$\phi$ (phase) & phase/vorticity ordering & skew-adjoint generator $R$ (energy-neutral) \\
$\chi$ (memory) & history dependence & positive-type kernel $K$ (dissipative) \\
\bottomrule
\end{tabular}
\end{center}
with exponential example $K_\tau(t)=\tfrac1\tau e^{-t/\tau}$, which reduces to an
internal variable. Three things follow:

\begin{enumerate}
\item \textbf{The phase layer is not $\theta=0$.} The $\theta=0$ of §\,3 is the
\emph{spatial} Gram metric, not HLV's time. HLV's phase sits in the skew-adjoint
operator $R$, dynamically driven by $\dot\phi(t)$. So the layer reads correctly:
\emph{dynamic operator-phase $R$} (HLV) vs \emph{fixed scalar phase $\theta=2/9$}
(FFGFT) -- not phaseless vs phased.

\item \textbf{The time/memory layer is a real bridge candidate.} HLV's $\chi$
channel is a positive-type, non-Markovian memory kernel $\int_0^t K(t-s)\,u(s)\,ds$.
FFGFT's time extension in Doc.~282 gives a memory kernel as the Fourier transform of
the discrete spectral density $\sum_k g_k^2\,\delta(\omega-\omega_k)$ -- likewise
non-Markovian (back-flow, revivals), with exponential form $\to$ internal variable.
\emph{The same class of object.} Where space forks (5- vs 3-fold), at the time level
two kernels of the same type face each other: \emph{a spatial fork, a possible
$\chi$-memory bridge}.

\item \textbf{Two honest limits.} First: the $i\!\cdot\!j$ product relation of the
spiral algebra is fixed in \emph{none} of Krüger's sources (Doc 271, FA-Hub,
Navier-Stokes) -- it is always bypassed via the operator translation. The phase
comparison must therefore be made at the \emph{operator} level, not the algebraic
one; whether operator-phase $R$ meets FFGFT's $2/9$ can only be decided with an
explicit $R$ specification. Second: Krüger himself records that many ordinary
reduced models generate memory; the HLV reading is only a candidate model class, not
a demonstrated ontology (Run~M1: the HLV kernel beats weak nulls, not the generic
biexponential). The $\chi$ bridge is thus \emph{structural kinship}, not a
uniqueness claim -- consistent with Doc.~276/281 ($\varphi$ is not specific either).
\end{enumerate}

\begin{important}[The three levels are not equally strong]
Reading the connection by physical level -- space, time, mass -- the three are
\emph{not} equivalent, and this is methodologically decisive:
\begin{itemize}
\item \textbf{Mass (Z$_3$ circulant):} the \emph{genuine} bridge -- derived from each
geometry on its own, with an explicit C$_3$ map and a real divergence point. It
alone fully passes the three-part criterion.
\item \textbf{Time ($\chi$ memory):} \emph{structural kinship}, a candidate -- not
HLV-unique (Run~M1: a generic biexponential beats the HLV kernel). A bridge
candidate, not a uniqueness claim.
\item \textbf{Space (connection-Laplace on a fibre bundle):} \emph{shared language},
not a bridge. That both are a connection-Laplace on a fibre-valued space holds for
\emph{every} lattice gauge theory -- this is Stufe-0 scaffolding, exactly what the
criterion excludes. It must not be counted as a load-bearing connection, or one
inflates the link toward triviality.
\end{itemize}
So only the mass level is a bridge in the sense of the criterion; the time level an
honest candidate; the space level a shared frame, no more.
\end{important}

% ==============================================================================
\section{What remains: hexagonal, not pentagonal}
% ==============================================================================

Strip HLV of the two axes along which it forks -- symmetry order and compactness --
and exactly one ingredient remains that is not already FFGFT: the spiral-time. Its
operators ($R$, $K$) depend on neither 3- nor 5-foldness and survive the switch to
3-fold unchanged. \emph{What substantially remains of HLV is therefore not a geometry
but a time} -- and it is there, in the dynamics, that it meets FFGFT's non-Markovian
kernel from Doc.~282. The geometric remainder -- icosahedral 5-foldness -- is an
elegant but physically empty branch: no sector realises it, and $\varphi$ is not
specific (Doc.~276/281).

A short geometric lesson follows. Five-foldness is \emph{non-crystallographic} -- it
does not tile periodically and forces decompactification; three-foldness is
\emph{crystallographic}, of the same family as six-foldness (the
triangular/hexagonal lattice). Where geometry carries the mass structure it is
therefore hexagonal, not pentagonal -- and this is no new finding but FFGFT's own
anchor: the \emph{Euler Tonnetz} is a hexagonal lattice (Doc.~275, ``the frame
anchors in the Euler Tonnetz, not in the golden ratio'').

\begin{important}[The defensible form of the thesis]
``Hexagonal'' means the crystallographic family (3-/6-fold), not the claim that
spacetime is a hexagonal tiling -- FFGFT's base remains $T^4$ with a $Z_3$ generation
symmetry. Empirically 3-fold-confirmed so far are only the charged leptons; the
hadrons are near-degenerate ($Q\approx1/3$), hence neither 3- nor 5-fold. What stands
is: \emph{where geometry carries the mass structure it is hexagonal/crystallographic;
the pentagonal stays an elegant, physically empty branch; and what substantially
remains of HLV is its time, not its geometry.}
\end{important}

% ==============================================================================
\section{Conclusion}
% ==============================================================================

All three parts of the criterion are answered. (1)~Same object: a Z$_3$ circulant
from HLV's own geometry. (2)~Map: the C$_3$-axis embedding of FFGFT's Z$_3$ into
HLV's icosahedral group. (3)~Divergence: fundamental, and in three layers --
compactness (decompactified/metric vs compact/$\tilde T\!\cdot\! m=1$), symmetry
order (5-fold $\varphi$, $1/\sqrt5$ vs 3-fold Koide, $\sqrt2$) and phase (HLV's
dynamic spiral-time operator-phase $R$ vs FFGFT's fixed scalar $2/9$; §\,8).
Compactness stays decisive for the \emph{mass} scale (no quantised modes, no
intrinsic $m$); HLV's \emph{time}, however, sits in the spiral-time, independently of
it. No parameter of HLV reaches FFGFT's protected Koide point -- yet at the
time/memory level both non-Markovian kernels are of the same class: \emph{a spatial
fork, a possible temporal bridge} (§\,8).

This is exactly what the criterion is built to deliver: a \emph{genuine} connection
-- shared Z$_3$-circulant family, explicit C$_3$ map -- \emph{with} a real
divergence point, and demonstrably \emph{not} Stufe-0 triviality (the protected
structure separates, only the meaningless part coincides). FFGFT and HLV sit at two
\emph{different} points of the same family. It is a \emph{fork}, not a merge: a
common stem (a Z$_3$ circulant from a higher-dimensional lattice with a hidden
space), two branches according to the symmetry order.

For the correspondence the statement is precise and fair: \emph{we meet in the
Z$_3$ circulant along your C$_3$ axis; we separate because your icosahedral 5-fold
symmetry rigidly forces $1/\sqrt5$, while FFGFT's Koide condition demands a 3-fold
$\sqrt2$ -- and no parameter of your construction bridges that.} Reproducible with
\texttt{ffgft\_hlv\_c3\_bridge\_probe.py} (six checks, numpy-only).
